\section{Exact recovery algorithms and their costs}\label{sec:enumeration}

\subsection{Binary affine enumeration}
Theorem~\ref{thm:general-moment} at $p=2$ gives the bounded moment
\begin{equation}\label{eq:binary-moment}
 \E 2^{r_2}=O_{a,b}(1),\qquad
 r_2=\dim_{\F_2}\ker(A-\delta I).
\end{equation}
For bits $z_i$ encode signs as $\tau_i=(-1)^{z_i}$. At vertex $i$,
the parity of the number of opposite-sign neighbors is
$\sum_{j\sim i}(z_i+z_j)=((A+DI)z)_i$ over $\F_2$.
Every desired assignment therefore solves
\begin{equation}\label{eq:binary-affine}
 (A+DI)z=(b\bmod2)\one\quad\text{over }\F_2.
\end{equation}
Since $D+\delta=2a$ is even, its homogeneous matrix is
$A-\delta I$ modulo two. The planted assignment guarantees consistency
in the promised model.

The algorithm first verifies that its input is a simple $D$-regular
graph, or reports failure, and solves \eqref{eq:binary-affine} by exact
binary elimination. If inconsistent, it reports failure. Otherwise
write the affine solution set as $z_0+\ker(A+DI)$.
If $2^{r_2}>n^{\nu+1}$, report failure. Enumerate the remaining affine
solutions in Gray-code order and, for each, count opposite-sign
neighbors over the integers. Return the first vector with exactly
$b$ such neighbors at every vertex, or failure if none passes.

Each accepted vector satisfies $A\tau=\delta\tau$ exactly and is
balanced. Checking parity alone would be insufficient: it permits
spurious modular candidates, especially at small characteristic.
By \eqref{eq:binary-moment} and Markov's inequality,
\begin{equation}\label{eq:cutoff-prob}
 \Prob\{2^{r_2}>n^{\nu+1}\}=O_{a,b}(N^{-\nu-1}).
\end{equation}
On the complementary event the candidate list contains the planted vector, and on
the rigidity event every accepted vector is $\pm\sigma$.
Theorem~\ref{thm:boolean} then proves the deterministic assertion
of Theorem~\ref{thm:main-recovery}, with total failure $O(N^{-\nu})$.

Binary elimination costs $O(n^3)$ elementary bit operations and
$O(n^2)$ storage using ordinary dense elimination. In Gray order one
basis vector is added per step, costing $O(n)$ bit operations;
the integer neighbor check costs $O_D(n)$ operations on bounded
integers. An $r_2$-bit counter also costs at most $O(n)$ per step.
Thus elimination and enumeration together cost
\begin{equation}\label{eq:binary-cost}
 O_D(n^3+n2^{r_2}),
\end{equation}
and the cutoff gives $O_{a,b}(n^3+n^{\nu+2})$ elementary work on every input.
This is a polynomial bit bound. In a sequential storage model the
addressing overhead is at most an additional polynomial factor; in
the usual bit-operation or word-RAM models the displayed bound, up
to logarithmic address costs, is immediate. There are no large
uncontrolled rational numbers in this algorithm.

Without the cutoff, the same procedure lists every passing assignment.
Suppress one representative of each global-sign pair, for example by
requiring its first sign to be positive. Equation~\eqref{eq:binary-moment}
and \eqref{eq:binary-cost} give expected work
$O_{a,b}(n^3)$, including output checks and a streamed output
list. The runtime can be exponential on rare inputs; the cutoff
version supplies the worst-case polynomial bound in the theorem.


A smaller universal cap $n^3$ gives bounded $O_D(n^4)$ elementary
work and failure $O(N^{-\nu})+O(N^{-3})$. In particular it gives
$O(N^{-2})$ recovery failure throughout the assortative range.
